\section{Discussion}
\label{sec:discussion}

%-------------------------------------------------------------------------
\subsection{Why the Original Diffusion Loss Fails}
\label{sec:diff_loss}
% [3 pts] Intuitively explain why L_Diff = ||eps_hat(x_t, y, t) - eps_t||^2 did not produce
% realistic samples (use < 1 math equation) & provide visual results of applying L_Diff.

To produce \cref{fig:diff_loss}, I used the same setup as SDS (CFG 25, lr 0.01, 500 steps) but minimized $\mathcal L_{\text{Diff}} = \|\hat\epsilon_\theta(x_t, y, t) - \epsilon\|^2$ directly, backpropagating through the UNet (separate script \texttt{scripts/diff\_loss.py}).
The results contain no recognizable object, only colorful high-frequency noise, while SDS with the same settings produces the expected content.

The difference is only one term.
The full gradient is $(\hat\epsilon_\theta - \epsilon)\,\frac{\partial \hat\epsilon_\theta}{\partial x_t}\,\frac{\partial x_t}{\partial x_0}$, and SDS simply drops the UNet Jacobian in the middle.
Intuitively, the diffusion loss does not ask the image to look real; it asks the UNet to predict the noise well.
Through the Jacobian, the optimizer finds that tiny high-frequency changes to the input are the fastest way to reduce the UNet's error, much like an adversarial attack, so the latent drifts towards patterns that are meaningless to humans.
Without the Jacobian, the residual $\hat\epsilon_\theta - \epsilon$ itself points towards a cleaner, prompt-consistent image, which is why SDS works.

%-------------------------------------------------------------------------
\subsection{Why SDI and VSD Outperform SDS}
\label{sec:why_better}
% [3 pts] Intuitively analyze why SDI and VSD perform better than SDS (without heavy math).

\noindent\textbf{SDS.}
The typical SDS artifacts in my results are blur, missing detail, and washed-out colors at normal guidance, and over-saturation at very high guidance (\cref{fig:qualitative,fig:abl_sds_gs}).
The cause is the random noise: every step uses a new $\epsilon$ and $t$, so each step pushes the image in a slightly different direction, and over many steps these directions average out into a smooth ``average'' image.
A large guidance scale is needed to keep a strong enough signal, but it also exaggerates colors.

\noindent\textbf{VSD.}
VSD replaces the random noise with a learned prediction of the current images.
The random part that SDS suffers from is now present on both sides of the difference and largely cancels, leaving a cleaner direction from ``what I generate now'' towards ``what the pretrained model expects''.
This is why VSD works with a normal guidance of 7.5 and gives much sharper details than SDS in \cref{fig:qualitative}.
In my runs it also introduced over-saturated textures, which lowered its CLIP score; I believe a better-tuned LoRA (\eg learning rate) would reduce them.

\noindent\textbf{SDI.}
SDI keeps the SDS update but chooses the noise by DDIM inversion instead of at random.
Because the noise is consistent with the current image, the target from one-step denoising is a sharper version of that image instead of a blurry average, and the latent converges smoothly towards it; the large-to-small timestep schedule further separates layout from detail.
This gives the cleanest images in \cref{fig:qualitative} at guidance 7.5.
The inversion ablation supports this: with a poor inversion (1 step) SDI collapses, while any reasonably accurate inversion gives good results.

Although SDS has the highest average CLIP score in my runs, the differences are small, and visually both VSD and SDI are clearly sharper and more detailed, which CLIP ViT-B/32 does not fully reflect.
